% Sections prepared independently for inclusion in the main manuscript.
% Requires amsmath, amssymb, and theorem/proposition/lemma/corollary/remark environments.
% No custom macros are used.

\section{Finite consequences: powers and minimal polynomials}
\label{sec:finite-applications}

Throughout this section, $n\geq1$ and
\[
 p(X)=X^n+\sum_{i=0}^{n-1}c_iX^i=\det(XI-A).
\]
All ring statements below are equational statements, or implications from
explicit finite data. To read them internally in intuitionistic primitive
recursive arithmetic, fix a primitive recursive presentation of the ring
whose ring laws have iPRA proofs. The field algorithms additionally require
primitive recursive equality and inversion of nonzero elements, with
iPRA proofs of the field axioms and the guarded inverse identities for
that fixed presentation. Rational
numbers and finite fields provide concrete examples. This convention does
not quantify over arbitrary rings as objects of first-order arithmetic.

\subsection{Monic division without a zero test}

\begin{lemma}[Fixed-length monic division]
For every coefficient list $f$ and monic $p$ of displayed degree $n$, there
are primitive recursively computed lists $q,r$ such that
\[
 f=qp+r,\qquad r_i=0\quad(i\geq n).
\]
They are unique as polynomials. No division by a ring element is used.
\end{lemma}
\begin{proof}
Use the input list length as an upper bound for degree. Starting at its last
index $d\geq n$, subtract $f_dX^{d-n}p$ and add $f_dX^{d-n}$ to the quotient.
The coefficient at $d$ becomes zero, regardless of whether it was already
zero. Continue down to $n$. The invariant is the coefficientwise identity
$f=q_{\mathrm{current}}p+r_{\mathrm{current}}$ with all processed high
coefficients zero. It is an equality of bounded lists and is preserved by
one distributive-law calculation.

For uniqueness, suppose $p(q-q')=r'-r$. Look at coefficients from the fixed
maximum list length downwards. The highest not-yet-eliminated coefficient
of the left side is the corresponding coefficient of $q-q'$, since the
leading coefficient of $p$ is $1$; the right side is zero in every degree
at least $n$. Descending induction therefore sets every coefficient of
$q-q'$ to zero, and then every coefficient of $r-r'$ to zero. This argument
needs neither an integral domain nor a decision of whether a coefficient
is zero.
\end{proof}

Evaluation in a matrix is multiplicative for scalar polynomials because
scalar matrices commute with $A$: expansion of the finite double sum gives
$(fg)(A)=f(A)g(A)$. Consequently $p(A)=0$ implies $f(A)=r(A)$.
For powers, one can avoid constructing the quotient. Let $r_{0}=1$ and
write $r_k=\sum_{i<n}r_{k,i}X^i$. Then
\begin{align*}
 r_{k+1,0}&=-c_0r_{k,n-1},\\
 r_{k+1,i}&=r_{k,i-1}-c_ir_{k,n-1}\qquad(1\leq i<n).
\end{align*}
The induction invariant is $A^k=\sum_{i<n}r_{k,i}A^i$.
Thus every nonnegative power has a computed representative of degree
less than $n$. Reduced multiplication followed by binary exponentiation
also computes this representative; binary exponentiation improves
efficiency but is unnecessary for primitive recursiveness.

Negative powers are also covered: use the explicit inverse from
Theorem~\ref{thm:inverse}, take its positive powers, and apply the same
monic reduction. The exact numerical examples were proved in
Section~\ref{sec:main}.

\subsection{Computing a minimal polynomial by bounded elimination}

\begin{proposition}[Effective field version]
Over a nontrivial discrete field with primitive recursive operations and
inversion, a primitive recursive algorithm constructs a monic polynomial
$m_A$ satisfying
\[
 m_A(A)=0,\qquad \deg m_A\leq n,\qquad
 h(A)=0\ \Longleftrightarrow\ m_A\mid h.
\]
In particular $m_A\mid p$; conversely that divisibility implies $p(A)=0$.
\end{proposition}
\begin{proof}
Flatten $I,A,A^2,\ldots$ into vectors $v_j$ of length $n^2$.
For $d=1,\ldots,n^2$, solve
\[
 [v_0\ \cdots\ v_{d-1}]x=-v_d
\]
by row reduction. Search for a pivot only within the remaining finite
rectangle. A nonzero pivot is inverted and its row is used to eliminate
its column. Every operation is invertible, so the transformed system is
equivalent to the original one. If a zero row of coefficients has nonzero
right-hand side, the stored product of row operations supplies a row
witness $w$ with $w[v_0\cdots v_{d-1}]=0$ but $wv_d\ne0$. Otherwise set
the free variables to zero and recover a solution by back substitution.
There is no search over possible field coefficients.

Some $d\leq n^2$ succeeds: row reduction of the $n^2\times(n^2+1)$
matrix $[v_0\cdots v_{n^2}]$ has at most $n^2$ pivots, so a free column
produces a dependence. Normalize its last nonzero coefficient. Since
$I\ne0$, that last index is positive. Select the first successful $d$
and its coefficients, obtaining $m_A=X^d+\sum_{j<d}x_jX^j$.
The supplied inconsistency witnesses for earlier $d$ certify minimality.
Cayley--Hamilton makes the system successful already at $d=n$, so the
output has $d\leq n$.

Divide any annihilator $h$ by $m_A$: $h=qm_A+r$, with $\deg r<d$.
Then $r(A)=0$. A bounded scan either verifies every coefficient of $r$
is zero, or finds its last nonzero coefficient. In the latter case,
division by that coefficient makes a monic annihilator of degree less
than $d$, contradicting the stored minimality certificates. A nonzero constant
remainder would imply $I=0$, which is excluded by the nontrivial field
and $n\ge1$ hypotheses. Thus $r=0$.
The converse follows from multiplicativity of evaluation. This also
proves uniqueness of the monic generator.
\end{proof}

The loops are bounded by displayed dimensions; their invariants are
equalities and decidable inequalities between finite arrays. A field
presented only by an abstract classical existence axiom for inverses
does not supply the algorithm assumed here. Nor is a minimal-polynomial
algorithm over an arbitrary commutative ring asserted.

\subsection{Companion matrices}

For a monic $f=X^n+a_{n-1}X^{n-1}+\cdots+a_0$, define its companion matrix
by
\[
 C_fe_i=e_{i+1}\ (i<n-1),\qquad
 C_fe_{n-1}=-\sum_{i<n}a_ie_i.
\]
Then $C_f^ie_0=e_i$ for $i<n$, so $f(C_f)e_0=0$. Commuting with
$C_f^i$ shows $f(C_f)e_i=0$ for every basis vector, and hence $f(C_f)=0$.
Expansion of $\det(XI-C_f)$ down the last column gives $f$: the cofactor
at row $i$ of that column is $X^i$. Indeed, its minor has $i$ diagonal
$X$ factors and $n-1-i$ subdiagonal $-1$ factors; its cofactor sign
cancels the latter signs. The final row contributes
$(X+a_{n-1})X^{n-1}$.

More strongly, over every commutative ring,
\[
 h(C_f)=0\quad\Longleftrightarrow\quad f\mid h.
\]
Divide $h=qf+r$ with $r$ supported below $n$. If $h(C_f)=0$, then
$0=r(C_f)e_0=\sum_{i<n}r_ie_i$, which says every $r_i$ is zero.
This argument does not require a field or linear dependence search.

\section{Eigenvalue statements with explicit finite data}

If $Av=\lambda v$, induction gives $A^kv=\lambda^kv$, and therefore
\[
 f(A)v=f(\lambda)v
\]
for every scalar polynomial $f$. With $f=p$, this yields $p(\lambda)v=0$.
Over a discrete field, a nonzero vector has a nonzero component found by
a bounded scan; inversion of that component gives $p(\lambda)=0$.
Over a ring the same conclusion holds if a component is supplied with a
unit inverse. Without such a hypothesis it can fail: in $\mathbb Z/6$,
$A=(2)$, $\lambda=5$, and $v=(2)$ satisfy $Av=\lambda v$ although
$p(5)=3\ne0$.

Conversely, over the effective field, $p(\lambda)=0$ means that
$\lambda I-A$ has zero determinant. Row reduction cannot have a pivot
in every column, because elementary row operations change determinant
by known nonzero factors and a triangular matrix with all pivots has
nonzero determinant. Set one free variable to $1$, the others to $0$,
and back substitute. The result is a certified nonzero eigenvector.
This proves the eigenvalue/root equivalence constructively under the
stated field assumptions.

Suppose matrices $S,T$ with $ST=TS=I$ are supplied and $B=TAS$.
Finite multiplication gives $B^k=TA^kS$, while determinant
multiplicativity gives
\[
 \det(XI-B)=\det T\,\det(XI-A)\,\det S=p(X).
\]
If $B$ is upper triangular with diagonal $\lambda_0,\ldots,\lambda_{n-1}$,
its determinant polynomial is $\prod_i(X-\lambda_i)$. To verify this
directly, a nonzero Leibniz summand would require $\sigma(i)\geq i$ for
every $i$; equality of the two finite sums of indices forces all
$\sigma(i)=i$. Products of triangular matrices multiply their diagonal
entries, so $\operatorname{tr}(B^k)=\sum_i\lambda_i^k$.
Finally $\operatorname{tr}(UV)=\operatorname{tr}(VU)$ follows by
interchanging two finite sums, and hence
\[
 \operatorname{tr}(A^k)=\sum_i\lambda_i^k.
\]
These are proofs from a supplied triangularization certificate; they do
not assume an algorithm finding roots or a Jordan form. A supplied
diagonalization is the special case in which $B$ is diagonal.

\section{Finitely generated modules and integral elements}

\subsection{The determinant trick without a basis}

\begin{proposition}[Finite generator form]
Let $g_0,\ldots,g_{n-1}$ be elements of an $R$-module, and let $\phi$ be
$R$-linear. Suppose
\[
 \phi(g_j)=\sum_{i<n}a_{ij}g_i.
\]
Writing $p=\det(XI-A)$, one has $p(\phi)(g_j)=0$ for every $j$, and
therefore $p(\phi)(\sum_jx_jg_j)=0$ for every supplied coordinate list
$x$. If these elements generate the module, $p(\phi)=0$ on it.
\end{proposition}
\begin{proof}
Let $\pi(x)=\sum_jx_jg_j$. The supplied equations say
$\phi\pi=\pi A$. Induction gives $\phi^k\pi=\pi A^k$; addition with the
coefficients of $p$ gives $p(\phi)\pi=\pi p(A)=0$. Apply to the unit
vectors, or to any supplied $x$. Linear independence of the generators
has never been used.
\end{proof}

In arithmetic the last assertion is read using a coordinate witness for
the element in question, or a primitive recursive coordinate function
when a uniform computational presentation is desired. It is not an
invocation of a choice principle to select coordinates in an arbitrary
abstract module.

If every $a_{ij}$ belongs to an ideal $J$, then the coefficient of
$X^{n-k}$ in $p$ belongs to $J^k$: the principal-minor formula displays
it as a finite sum of products of $k$ such entries. This proves the
usual strengthened determinant trick with explicit membership
certificates.

For a certificate version of Nakayama's lemma, suppose
$g_j=\sum_i a_{ij}g_i$ with $a_{ij}\in J$. Apply the proposition with
$\phi=1$. The scalar
$\delta=\det(I-A)=1-b$, for a displayed $b\in J$, annihilates every
$g_j$. If an inverse $u\delta=1$ is supplied, multiplying gives $g_j=0$.
Thus the whole generated module is zero. The familiar assumption that
$J$ lies in the Jacobson radical supplies this unit property in the
abstract formulation; the finite arithmetic theorem uses the inverse
as data and does not search for it.

\subsection{Closure of integral elements with a degree bound}

\begin{proposition}[Explicit integrality construction]
In a commutative $R$-algebra, suppose $\alpha_j$ satisfies a supplied
monic polynomial $f_j$ of positive degree $d_j$, for $1\leq j\leq k$.
For every supplied polynomial expression $\beta$ in the $\alpha_j$,
one can construct a monic annihilator of $\beta$ of degree
$D=\prod_{j=1}^k d_j$.
\end{proposition}
\begin{proof}
Use the $D$ monomials
$g_e=\prod_j\alpha_j^{e_j}$, $0\leq e_j<d_j$, including $g_0=1$.
Reduce high powers by the supplied monic relations, one variable at a
time. Every reduction preserves the represented element and decreases
the exponent currently being processed. It computes coefficients for
$\beta g_e$ in the displayed list. Let $A$ be the resulting $D\times D$
matrix. Multiplication by $\beta$ satisfies the generator equations,
so the determinant trick gives $p_A(\beta)g_0=0$. As $g_0=1$,
$p_A(\beta)=0$. The construction uses lists and bounded monic division.
It neither assumes the monomials are independent nor computes a basis
of their span.
\end{proof}

Taking $R=\mathbb Z$ proves closure of integral elements under addition,
multiplication, and negation, with a computed monic equation. In
particular two input equations of degrees $d,e$ give an equation of
degree $de$ for both the sum and the product. This bypasses any theorem
that a submodule of a finitely generated module is finitely generated.

For an element $\alpha$ of a finite extension presented by a
$\mathbb Q$-basis, a unit coordinate vector, and a rational multiplication
table, form the matrix $L_\alpha$ of multiplication by $\alpha$.
Associativity and the unit give
$h(L_\alpha)=L_{h(\alpha)}$ and
$L_{h(\alpha)}(1)=h(\alpha)$. Thus its matrix minimal polynomial is
exactly the minimal polynomial of $\alpha$. The bounded algorithm above
computes it. Computing the characteristic polynomial alone only gives
an annihilator; it need not be minimal. For example the multiplication
matrix of $1$ in any extension of degree $D>1$ has characteristic
polynomial $(X-1)^D$ and minimal polynomial $X-1$.

If $\alpha$ is integral over $\mathbb Z$, its rational minimal polynomial
has integer coefficients. Here is the finite algebra needed for that
extra claim. Let $P\in\mathbb Z[X]$ be a monic annihilator and
$P=mq$ in $\mathbb Q[X]$, with $m,q$ monic. Clear denominators and
divide coefficient gcds to write $f=am$, $g=bq$ as primitive integer
polynomials with positive leading coefficients $a,b$. These are positive
integers: the coefficient gcd divides the leading coefficient of the
cleared monic polynomial. A product of
primitive integer polynomials is primitive: otherwise a prime dividing
all its coefficients, found by bounded search from their gcd, would
make its reduction zero modulo that prime. The two factors each have
a last nonzero reduced coefficient; their product gives a nonzero
coefficient in the sum of these two degrees, a contradiction. Hence
$fg=abP$ has content $1$, whereas the right side has content $ab$.
Therefore $a=b=1$ and $m\in\mathbb Z[X]$. Integer gcd, bounded prime
search, and finite-field multiplication all have primitive recursive
implementations.

\section{Interpolation from supplied splitting data}

Suppose a monic annihilating polynomial is supplied in the form
\[
 p(X)=\prod_{j<r}(X-\lambda_j)^{m_j},\qquad
 m_j\geq1,\quad\sum_jm_j=n,
\]
and every difference $\lambda_i-\lambda_j$ for $i\ne j$ has a supplied
unit inverse. These assumptions are stronger than the mere existence of
roots and make all interpolation operations explicit.

\subsection{Distinct roots and Sylvester's formula}

If all $m_j=1$, set
\[
 L_j(X)=\prod_{i\ne j}\frac{X-\lambda_i}{\lambda_j-\lambda_i}.
\]
For every polynomial $f$,
\[
 \operatorname{rem}_p f=\sum_j f(\lambda_j)L_j(X),\qquad
 f(A)=\sum_jf(\lambda_j)L_j(A).
\]
To prove uniqueness, a polynomial vanishing at $\lambda_0$ is divisible
by $X-\lambda_0$, by monic division. If it also vanishes at $\lambda_1$,
evaluate the quotient identity there and cancel the supplied unit
$\lambda_1-\lambda_0$ to obtain the next factor. Continuing shows that
a polynomial of degree less than $n$ vanishing at every root is zero.
The displayed candidate has the required values, so it is the
remainder. Evaluating its congruence modulo $p$ gives the matrix formula.
The same argument applied to products shows
$L_i(A)L_j(A)=0$ for $i\ne j$, $L_j(A)^2=L_j(A)$, and
$\sum_jL_j(A)=I$.

\subsection{Repeated roots: jets and Hasse derivatives}

The $s$th Hasse derivative at $\lambda$ is the coefficient of $Z^s$ in
$f(\lambda+Z)$:
\[
 D^{[s]}f(\lambda)=\sum_{k\geq s}\binom{k}{s}f_k\lambda^{k-s}.
\]
This is a finite formula in every characteristic. Ordinary derivatives
satisfy $f^{(s)}(\lambda)=s!D^{[s]}f(\lambda)$, so ordinary derivative
conditions are equivalent only when the relevant factorials are units.

Let $q_j=p/(X-\lambda_j)^{m_j}$. Write
$q_j(\lambda_j+Z)=\sum_s a_{j,s}Z^s$ and put
\[
 b_{j,0}=a_{j,0}^{-1},\qquad
 b_{j,t}=-a_{j,0}^{-1}\sum_{s=1}^t a_{j,s}b_{j,t-s}
 \quad(1\leq t<m_j).
\]
The constant $a_{j,0}=\prod_{i\ne j}(\lambda_j-\lambda_i)^{m_i}$
has a supplied inverse. Convolution proves
$q_j(\lambda_j+Z)\sum_{t<m_j}b_{j,t}Z^t\equiv1\pmod{Z^{m_j}}$.
Thus
\[
 E_j(X)=q_j(X)\sum_{t<m_j}b_{j,t}(X-\lambda_j)^t
\]
is $1$ modulo $(X-\lambda_j)^{m_j}$ and $0$ modulo every other factor.
For arbitrary jet data $d_{j,s}$ the polynomial
\[
 r=\operatorname{rem}_p\left(
 \sum_jE_j(X)\sum_{s<m_j}d_{j,s}(X-\lambda_j)^s\right)
\]
therefore has exactly these jets and degree less than $n$.

For uniqueness, note that the same truncated inverse construction shows
that any product of the other factors is invertible modulo
$(X-\lambda_j)^{m_j}$. If a polynomial is divisible by that product
and by $(X-\lambda_j)^{m_j}$, its quotient by the former is divisible by
the latter: multiply its zero residue by this explicit inverse.
Induction therefore makes a polynomial divisible by each factor
divisible by their product $p$. A difference of two candidates of
degree less than $n$ must vanish by monic division uniqueness.
Taking $d_{j,s}=D^{[s]}f(\lambda_j)$ proves $f(A)=r(A)$.

This explains why repeated eigenvalues require derivative data in
characteristic zero, and why repeated values alone do not determine a
unique remainder. The proof also supplies a division-free replacement
for derivatives in positive characteristic, apart from the explicitly
assumed inverse root differences.

\section{Matrix functions as controlled finite calculations}

An arbitrary analytic function is not a first-order arithmetic object.
Here the precise replacement is a specified primitive recursive sequence
of rational coefficients and a primitive recursive modulus for finite
tails. Equalities of the corresponding Cauchy names mean explicit
finite inequalities at every precision. The external real-number
interpretation is useful notation, not an additional completeness axiom
inside iPRA.

\subsection{A sufficient quantitative polynomial-reduction theorem}

Let $A$ be a rational matrix and let $r_k=\operatorname{rem}_pX^k$ as
above. For a polynomial supported below $n$, write
$|r|_1=\sum_{i<n}|r_i|$, and choose rational
\[
 K\geq\max\left(1,\sum_{i<n}|c_i|\right),\qquad
 H\geq\max(1,\|A\|_\infty^{n-1}),
\]
where $\|A\|_\infty=\max_i\sum_j|a_{ij}|$.
The recurrence for $r_k$ gives
$|r_{k+1}|_1\leq K|r_k|_1$, so $|r_k|_1\leq K^k$.
The finite triangle inequality and submultiplicativity of the row norm
give $\|r(A)\|_\infty\leq H|r|_1$.

Suppose a primitive recursive rational sequence $a_k$ and a primitive
recursive function $N$ are supplied together with proofs that
\[
 \sum_{k=N(s)}^m |a_k|K^k\leq2^{-s}
 \qquad(m\geq N(s)).                                      \tag{*}
\]
Define finite approximants
\[
 \beta_{i,M}=\sum_{k<M}a_kr_{k,i},\qquad
 F_M=\sum_{k<M}a_kA^k.
\]
Their exact finite identity is
$F_M=\sum_{i<n}\beta_{i,M}A^i$. Condition (*) bounds every coefficient
tail by $2^{-s}$ and every matrix tail by $H2^{-s}$. It follows, in the
precise sense of these finite estimates, that
\[
 f(A)=\sum_{i<n}\beta_iA^i,\qquad
 \beta_i=\sum_{k\geq0}a_kr_{k,i}.
\]
For a requested matrix error $2^{-e}$, choose an integer $h$ with
$H\leq2^h$ and use $N(e+h)$. Thus the conclusion contains an explicit
primitive recursive modulus. It does not assert that every analytic
function near the spectrum satisfies (*) for this possibly larger $K$.

The statement $f=qp+r$ with both $q$ and $r$ polynomials cannot hold
for a nonpolynomial analytic $f$. Even when an analytic quotient is
available, that is a different assertion. Our argument uses only
polynomial identities for finite partial sums and the specified tail
bounds.

\subsection{An explicit factorial tail bound}

For a positive integer $U$ and $s\geq0$, put
\[
 M(U,s)=2U+2U^2+s+2.
\]
For $k\geq2U$, the ratio of successive terms $U^k/k!$ is at most
$1/2$. Also $U^{2U}/(2U)!\leq U^{2U}\leq2^{2U^2}$, since
$U\leq2^U$. Therefore
\[
 \frac{U^{M(U,s)}}{M(U,s)!}\leq2^{-s-2},\qquad
 \sum_{k=M(U,s)}^v\frac{U^k}{k!}\leq2^{-s-1}
 \quad(v\geq M(U,s)).
\]
The last inequality is the finite geometric-sum bound. All denominators
can be cleared, leaving inequalities of natural numbers proved by
bounded summation and quantifier-free induction.

For rational $t$, choose an integer $U\geq\max(1,|t|K)$.
The exponential coefficients $a_k=t^k/k!$ satisfy (*) with this $M$;
the sine and cosine coefficient subsequences satisfy the same bound.
If a direct matrix-series bound is also wanted, enlarge $U$ to exceed
$|t|\|A\|_\infty$. Consequently matrix exponential, sine, and cosine
have primitive recursive rational Cauchy names and polynomial
representations of degree less than $n$, with the displayed moduli.

\subsection{The article's examples and rotation identities}

For $T=\left(\begin{smallmatrix}1&2\\0&3\end{smallmatrix}\right)$,
induction gives
\[
 T^k=\begin{pmatrix}1&3^k-1\\0&3^k\end{pmatrix}\qquad(k\geq0).
\]
Hence every polynomial obeys
\[
 f(T)=\begin{pmatrix}f(1)&f(3)-f(1)\\0&f(3)\end{pmatrix}.
\]
Apply this exact identity to exponential or sine partial sums and use
the factorial estimate. For every rational $t$, as identities of the
specified Cauchy names,
\[
 e^{tT}=\begin{pmatrix}e^t&e^{3t}-e^t\\0&e^{3t}\end{pmatrix},\qquad
 \sin(tT)=\begin{pmatrix}\sin t&\sin3t-\sin t\\0&\sin3t\end{pmatrix}.
\]
The corresponding coefficients in $c_0I+c_1T$ are
$c_0=(3f(1)-f(3))/2$ and $c_1=(f(3)-f(1))/2$.

More generally the relation $J^2=-I$ gives
$J^{2r}=(-1)^rI$ and $J^{2r+1}=(-1)^rJ$, whence
\[
 e^{tJ}=(\cos t)I+(\sin t)J.
\]
For $J=\left(\begin{smallmatrix}0&1\\-1&0\end{smallmatrix}\right)$
this is the article's two-dimensional rotation matrix. If $P^2=I$,
the same calculation over rational complex numbers gives
$e^{itP}=(\cos t)I+i(\sin t)P$. In particular, with the usual Pauli
matrices
\[
 \sigma_1=\begin{pmatrix}0&1\\1&0\end{pmatrix},\qquad
 \sigma_2=\begin{pmatrix}0&-i\\i&0\end{pmatrix},\qquad
 \sigma_3=\begin{pmatrix}1&0\\0&-1\end{pmatrix},
\]
direct multiplication gives
$\sigma_j^2=I$ and
$\sigma_j\sigma_k+\sigma_k\sigma_j=0$ for $j\ne k$.
Thus $P=\sum_j u_j\sigma_j$ has $P^2=I$ whenever the supplied rational
coordinates satisfy $\sum_j u_j^2=1$.

For rational complex matrices, use the scalar norm
$|a+bi|_1=|a|+|b|$ and its induced row-sum matrix norm. Direct expansion
proves submultiplicativity, so the same factorial estimates justify
the complex-valued formulas above.

For a rational unit vector $u=(u_1,u_2,u_3)$ put
\[
 K=\begin{pmatrix}0&-u_3&u_2\\u_3&0&-u_1\\-u_2&u_1&0\end{pmatrix}.
\]
Entrywise multiplication gives $K^2=uu^{\mathsf T}-I$ and $Ku=0$,
so $K^3=-K$. Hence $K^{2r+1}=(-1)^rK$ and
$K^{2r+2}=(-1)^rK^2$, yielding Rodrigues' formula
\[
 e^{tK}=I+(\sin t)K+(1-\cos t)K^2.
\]
Each displayed analytic identity is obtained from a finite identity
between partial sums followed by the same explicit tail estimates.
Here $\cos t$ and $\sin t$ denote the Cauchy names given respectively by
the finite sums of $(-1)^rt^{2r}/(2r)!$ and
$(-1)^rt^{2r+1}/(2r+1)!$; the exponential uses $t^r/r!$.
No extraction of eigenvalues, analytic continuation, or appeal to a
matrix Lie group is involved. Rational parameters give a fully
arithmetical family; extending it to further real inputs requires
their approximation codes and error bounds as additional data.

\section{Quaternion and octonion boundaries}

\subsection{A precise quaternionic Cayley--Hamilton theorem}

For a rational quaternion $q=a+bi+cj+dk$ define
\[
 L(q)=\begin{pmatrix}
 a&-b&-c&-d\\
 b&a&-d&c\\
 c&d&a&-b\\
 d&-c&b&a
 \end{pmatrix}.
\]
The columns are the coordinates of $q$, $qi$, $qj$, and $qk$.
The quaternion multiplication table proves
$L(qr)=L(q)L(r)$, $L(q+r)=L(q)+L(r)$, and $L(1)=I_4$.
The first column recovers $q$, so $L$ is injective.
Apply $L$ blockwise to $A\in M_n(\mathbb H_{\mathbb Q})$, producing
$L_n(A)\in M_{4n}(\mathbb Q)$. Matrix multiplication is preserved because
each block product is a finite sum of products preserved by $L$.
If $p=\det(XI-L_n(A))$, then
\[
 L_n(p(A))=p(L_n(A))=0.
\]
Injectivity gives $p(A)=0$. This is an explicit monic rational
annihilator of degree $4n$, not a naive quaternionic Leibniz determinant.

There is a degree $2n$ refinement. In $\mathbb Q(i)$ set
\[
 \rho(q)=\begin{pmatrix}a+bi&c+di\\-c+di&a-bi\end{pmatrix},\qquad
 J=\begin{pmatrix}0&1\\-1&0\end{pmatrix}.
\]
Multiplication of the four basis matrices verifies the quaternion
relations, so $\rho$ is an injective algebra homomorphism. Furthermore
$J\overline{\rho(q)}J^{-1}=\rho(q)$. For the block matrix
$B=\rho_n(A)$ and block diagonal $J_n$, this gives
$J_n\overline B J_n^{-1}=B$. Determinant similarity therefore makes
every coefficient of $\det(XI-B)$ fixed by complex conjugation.
A rational complex number fixed by conjugation has zero imaginary
coordinate, so this characteristic polynomial lies in $\mathbb Q[X]$.
Its degree is $2n$ and its evaluation at $A$ is zero by injectivity.
All representation identities involve finitely many rational coordinate
equalities; no quaternionic spectral theorem is required.

For split quaternions with $i^2=j^2=1$ and $ij=k=-ji$, the representation
\[
 a+bi+cj+dk\longmapsto
 \begin{pmatrix}a+b&c+d\\c-d&a-b\end{pmatrix}
\]
is faithful over $\mathbb Q$ by its four coordinate equations.
Blockwise application gives a monic degree $2n$ rational annihilator in
the same way. These are carefully specified instances of the article's
noncommutative generalizations.

\subsection{What the octonion assertion does and does not give}

For one rational octonion $q=a_0+\sum_{i=1}^7a_ie_i$, direct use of the
multiplication table gives
$q\overline q=\sum_{i=0}^7a_i^2$ and $\overline q=2a_0-q$. Distributing
the product, which involves only two factors, proves
\[
 q^2-2a_0q+\left(\sum_{i=0}^7a_i^2\right)1=0.
\]
This is an elementary quadratic identity. It is not a proof of the
modified theorem for matrices over octonions mentioned in the article.
The latter requires a specified nonassociative evaluation convention
and a different theorem, and is not claimed here.

\subsection{Arithmetic status of these applications}

The constructions above use only fixed primitive recursive ring
operations, coefficient lists, products, bounded searches, and
primitive recursion with displayed bounds or parameters. Their stated
invariants are finite arrays of equalities and decidable inequalities.
Finite conjunctions are represented by a Boolean fold, so induction on
such invariants is quantifier-free induction after expansion of the
primitive recursive definitions. Algebraic structure axioms are used
as supplied equations; field inverses, root differences, generator
coordinates, and analytic tail estimates are additional data precisely
where indicated. This gives an explicit route to iPRA proofs of these
arithmetized statements. It does not turn an abstract host-language
proof about all rings, all modules, or all real functions into an
object-language iPRA derivation without the required coding and proof
translation.
